
% ============================================================

% --- Title Block ---
\begin{center}
  {\LARGE\bfseries\color{deepblue}
    Stochastic Pursuit on a Dynamic Torus}\\[0.6em]
  {\large\color{midblue}
    A Composable Geometric Simulation Engine for Stochastic Pursuit,\\[0.2em]
    Stealth, Pack Strategy, and Emergent Neural Behaviour\\[0.2em]
    on Dynamic Riemannian Manifolds, Rendered via Vulkan}\\[1.2em]
  {\normalsize\color{softgold}\itshape
    Research \& Study Notes\quad$\cdot$\quad Solo Independent Study}\\[0.4em]
  {\small\color{gray} \today}
\end{center}

\vspace{1.2em}
\noindent\rule{\linewidth}{0.8pt}
\vspace{0.6em}

% --- Torus Visualization ---
\tdplotsetmaincoords{65}{125}
\begin{center}
\begin{tikzpicture}[tdplot_main_coords, scale=1.15]

  \def\R{2.2}
  \def\r{0.85}
  \colorlet{meshcol}{blue!30!gray}
  \colorlet{surfcol}{blue!7!white}

  % Back parallels
  \foreach \vv in {0,40,80,120,160,200,240,280,320} {
    \draw[meshcol!45, very thin]
      plot[domain=100:260, samples=28, variable=\u]
      ({(\R + \r*cos(\vv))*cos(\u)},
       {(\R + \r*cos(\vv))*sin(\u)},
       {\r*sin(\vv)});
  }
  % Back meridians
  \foreach \uu in {120,160,200,240} {
    \draw[meshcol!45, very thin]
      plot[domain=0:360, samples=36, variable=\v]
      ({(\R + \r*cos(\v))*cos(\uu)},
       {(\R + \r*cos(\v))*sin(\uu)},
       {\r*sin(\v)});
  }

  % Surface fill patches (front half)
  \foreach \vv in {0,20,...,340} {
    \pgfmathsetmacro{\vnext}{\vv+20}
    \fill[surfcol, opacity=0.6]
      plot[domain=-82:100, samples=18, variable=\u]
      ({(\R + \r*cos(\vv))*cos(\u)},
       {(\R + \r*cos(\vv))*sin(\u)},
       {\r*sin(\vv)})
      --
      plot[domain=100:-82, samples=18, variable=\u]
      ({(\R + \r*cos(\vnext))*cos(\u)},
       {(\R + \r*cos(\vnext))*sin(\u)},
       {\r*sin(\vnext)})
      -- cycle;
  }

  % Front parallels
  \foreach \vv in {0,40,80,120,160,200,240,280,320} {
    \draw[meshcol, thin]
      plot[domain=-82:100, samples=28, variable=\u]
      ({(\R + \r*cos(\vv))*cos(\u)},
       {(\R + \r*cos(\vv))*sin(\u)},
       {\r*sin(\vv)});
  }
  % Front meridians
  \foreach \uu in {-80,-40,0,40,80} {
    \draw[meshcol, thin]
      plot[domain=0:360, samples=36, variable=\v]
      ({(\R + \r*cos(\v))*cos(\uu)},
       {(\R + \r*cos(\v))*sin(\uu)},
       {\r*sin(\v)});
  }

  % Outer equator ring
  \draw[gray!35, thin]
    plot[domain=0:360, samples=72, variable=\u]
    ({(\R+\r)*cos(\u)},{(\R+\r)*sin(\u)},{0});

  % LEADER: orange-red, u=28, v=18
  \pgfmathsetmacro{\Lx}{(\R + \r*cos(18))*cos(28)}
  \pgfmathsetmacro{\Ly}{(\R + \r*cos(18))*sin(28)}
  \pgfmathsetmacro{\Lz}{\r*sin(18)}
  \node[circle, fill=orange!95!red, draw=red!70!black,
        inner sep=3.5pt, line width=0.7pt]
    at (\Lx,\Ly,\Lz) {};
  \node[font=\footnotesize\bfseries, text=orange!80!red,
        above right=-1pt and 3pt]
    at (\Lx,\Ly,\Lz) {Leader};

  % CHASER 1: cyan, u=55, v=348
  \pgfmathsetmacro{\CAx}{(\R + \r*cos(348))*cos(55)}
  \pgfmathsetmacro{\CAy}{(\R + \r*cos(348))*sin(55)}
  \pgfmathsetmacro{\CAz}{\r*sin(348)}
  \node[circle, fill=cyan!85!blue, draw=blue!70!black,
        inner sep=3.0pt, line width=0.7pt]
    at (\CAx,\CAy,\CAz) {};
  \node[font=\footnotesize\bfseries, text=cyan!55!blue,
        below right=0pt and 3pt]
    at (\CAx,\CAy,\CAz) {Chaser 1};

  % CHASER 2: magenta, u=8, v=42
  \pgfmathsetmacro{\CBx}{(\R + \r*cos(42))*cos(8)}
  \pgfmathsetmacro{\CBy}{(\R + \r*cos(42))*sin(8)}
  \pgfmathsetmacro{\CBz}{\r*sin(42)}
  \node[circle, fill=magenta!80!red, draw=magenta!60!black,
        inner sep=3.0pt, line width=0.7pt]
    at (\CBx,\CBy,\CBz) {};
  \node[font=\footnotesize\bfseries, text=magenta!65!black,
        above left=0pt and 3pt]
    at (\CBx,\CBy,\CBz) {Chaser 2};

  % Dashed pursuit lines
  \draw[cyan!65, dashed, thick, opacity=0.85]
    (\CAx,\CAy,\CAz) -- (\Lx,\Ly,\Lz);
  \draw[magenta!65, dashed, thick, opacity=0.85]
    (\CBx,\CBy,\CBz) -- (\Lx,\Ly,\Lz);

\end{tikzpicture}

\vspace{-0.3em}
{\small\itshape\color{gray}
  $\mathbb{T}^2 = \mathbb{R}^2/\mathbb{Z}^2$ with mesh,
  leader~({\color{orange!80!red}$\bullet$}) and
  chasers~({\color{cyan!60!blue}$\bullet$}\ {\color{magenta!65!black}$\bullet$}).
  Dashed lines indicate pursuit geodesics.}
\end{center}

\vspace{0.3em}
\noindent\rule{\linewidth}{0.4pt}
\vspace{0.3em}

% ============================================================
\section*{Mission Statement}
% ============================================================

This project is a long-horizon study in the unity of pure mathematics and
high-performance computation. The goal is to build, from first principles, a
\textbf{composable geometric simulation engine} capable of hosting arbitrary
families of stochastic agents on dynamic Riemannian manifolds, rendered in
real time via Vulkan.

The engine is not a single simulation but a \emph{laboratory}. Agents---
predators, prey, Brownian walkers, stealthy hunters, coordinated packs,
neural-network-driven evolvers---are declared by implementing a common
geometric interface. Effects on the arena---metric ripples triggered by user
pokes, potential wells, damping fields, geodesic flows---are independent
composable components. Every combination is geometrically correct by
construction, because the interface enforces it.

The mathematics dictates the architecture. Every ``thing'' in the
simulation is a \emph{section of a bundle} over the manifold: drift is a
tangent vector field, force is a cotangent field, the metric perturbation is
a symmetric $(0,2)$-tensor field, noise is a bundle morphism into the
tangent space. Composability is fiberwise addition and composition of bundle
maps---not a software pattern imposed on the mathematics, but the mathematics
itself transcribed into a type system.

The simulation is not an end. It is a \emph{forcing function}---a concrete
artifact that demands genuine understanding at every layer. A renderer that
produces a wrong geodesic is a proof of misunderstanding. A stealthy pursuer
that violates its curvature budget is a failed constraint. A pack that cannot
topologically trap its prey has not understood the fundamental group of the
torus. The Vulkan pipeline is an automated grader with no partial credit.

The work proceeds without shortcuts and without a deadline. Each stage of the
dependency graph below must be structurally understood before the next is
approached. This document is the living front matter of that journey.

% ============================================================
\section*{Problem Statement}
% ============================================================

\subsection*{The Arena}

Let $\mathbb{T}^2$ denote the flat two-dimensional torus, realized as the
quotient:
\[
  \mathbb{T}^2 \;\cong\; \mathbb{R}^2 / \mathbb{Z}^2,
\]
equipped with the background Riemannian metric $g_0 = dx^2 + dy^2$. This
metric is flat (zero Gaussian curvature) and induces a geodesic distance
\[
  d_0(x, y) \;=\; \min_{(m,n)\in\mathbb{Z}^2}
    \left\| (x - y) + (m,n) \right\|_{\mathbb{R}^2}.
\]
The arena is compact, boundary-free, and doubly periodic---a natural domain
for stochastic processes that must neither escape to infinity nor reflect from walls.

\subsection*{The Dynamic Metric}

At any time $t \geq 0$ and in response to user input (a ``poke'' at point
$p \in \mathbb{T}^2$ initiated at time $t_0$), the metric is deformed
conformally:
\[
  g(x, t) \;=\; \bigl(1 + \varepsilon\, h(x, t;\, p, t_0)\bigr)\, g_0,
  \qquad \varepsilon A < 1,
\]
where $h$ is a smooth, spatially propagating, temporally decaying perturbation
field. The positive-definiteness condition $\varepsilon A < 1$ guarantees that
$g(x,t)$ remains a valid Riemannian metric for all time. The perturbed geodesic
distance satisfies, to first order in $\varepsilon$:
\[
  d_g(x, y, t) \;=\; d_0(x,y)
    \;+\; \frac{\varepsilon}{2}\int_\gamma h(\gamma(s), t)\,ds
    \;+\; O(\varepsilon^2),
\]
where $\gamma$ is the geodesic in the flat metric. As the perturbation decays,
$g(x,t) \to g_0$ smoothly and the arena returns to flatness.

\subsection*{Stealth via Geodesic Curvature Constraint}

A na\"ive pursuer travels the shortest geodesic toward its target---maximally
efficient but maximally detectable. A \emph{stealthy} pursuer minimises a
detectability functional while still closing distance. The natural geometric
formalisation penalises the \textbf{geodesic curvature} of the pursuer's path:
\[
  \kappa_g(t) \;=\; \left\| \frac{D\dot{\gamma}}{dt} \right\|_{g(t)},
\]
where $D/dt$ denotes the covariant derivative along $\gamma$ with respect to
the instantaneous metric $g(t)$. Stealth is imposed as a kinematic constraint:
\[
  \kappa_g(t) \;\leq\; \kappa_{\max} \quad \forall\, t.
\]
The stealthy pursuit problem is then a constrained optimal control problem on
a Riemannian manifold---an isoperimetric problem in the sense of the calculus
of variations, with the curvature bound as the isoperimetric condition.

\subsection*{Pack Strategy and Topological Trapping}

A single pursuer on $\mathbb{T}^2$ cannot simultaneously block both
non-contractible escape routes. The fundamental group
$\pi_1(\mathbb{T}^2) = \mathbb{Z} \times \mathbb{Z}$ has two independent
generators---two families of loops that cannot be contracted to a point.
A \emph{pack} of $n$ pursuers coordinates to eliminate these escape routes
through a shared strategy encoded as a vector field on the configuration
manifold $(\mathbb{T}^2)^n$.

The pack's collective state is measured by the \textbf{coverage functional}:
\[
  \Phi(t) \;=\; \mathrm{Vol}\!\left(
    \mathbb{T}^2 \setminus \bigcup_{i=1}^{n}
    B_{g(t)}\!\left(X_i(t),\, r_i\right)
  \right),
\]
where $B_{g(t)}(X_i, r_i)$ is the geodesic ball of radius $r_i$ in the
instantaneous metric. The pack wins when $\Phi(t) = 0$. The minimum pack
size $n^*$ required to guarantee $\Phi \to 0$ is a topological invariant
of the arena, determined by the homology of $\mathbb{T}^2$.

\subsection*{Neural Policy Evolution}

Beyond analytically prescribed strategies, this project studies \emph{learned}
pursuit and evasion. Each agent's strategy is a smooth policy map:
\[
  \pi_\theta \;:\; \mathcal{S} \;\longrightarrow\; T\mathbb{T}^2,
\]
from a geometric state space $\mathcal{S}$ (encoding positions, delayed
observations, local metric, and curvature) to the tangent bundle---the space
of admissible velocities. The network must respect the periodicity of
$\mathbb{T}^2$ and any active constraints.

The evolutionary pressure operates on a \emph{population} of policies,
selecting jointly for capture efficiency, stealth score, and coverage
of the torus. This is a gradient flow on the \textbf{Wasserstein manifold}
of probability measures over the policy space---a research frontier in
multi-agent reinforcement learning and optimal transport.

\subsection*{The Agent SDE}

Let $\{X_i(t)\}_{i=1}^{n}$ be a collection of stochastic processes on
$\mathbb{T}^2$. Each agent satisfies a stochastic differential equation of the
form:
\[
  dX_i^k(t) \;=\; V_i^k\!\bigl(X_i(t),\, X_j(t - \tau)\bigr)\,dt
  \;+\; \sigma^k_{\ \ell}(X_i(t), t)\,dW_i^\ell(t),
\]
where:
\begin{itemize}
  \item $V_i$ is the \textbf{pursuit drift vector field}, steering agent $i$
    toward another agent's position at the delayed time $t - \tau$;
  \item $\tau > 0$ is the \textbf{information delay}---the pursuer reacts to
    where the prey \emph{was}, not where it \emph{is};
  \item $\sigma \sigma^T = g^{-1}(t)$ encodes the geometry of the arena into
    the noise covariance, so that Brownian fluctuations are isotropic \emph{in
    the metric sense} at each moment;
  \item $W_i$ are independent standard Wiener processes.
\end{itemize}

The generator of this process is the time-dependent operator:
\[
  \mathcal{L}_t \;=\; \frac{1}{2}\Delta_{g(t)} \;+\; V_i \cdot \nabla,
\]
where $\Delta_{g(t)}$ is the \textbf{Laplace-Beltrami operator} of the
instantaneous metric---the object that connects stochastic analysis,
differential geometry, and the heat equation in a single symbol.

\subsection*{The Central Questions}

This project asks, and will eventually answer, the following:

\begin{enumerate}
  \item \textbf{Well-posedness.} Under what conditions on $\tau$, $\varepsilon$,
    and $\sigma$ does the system of SDEs admit a unique strong solution on
    $\mathbb{T}^2$?

  \item \textbf{Geodesic bifurcation.} When a metric ripple transiently inflates
    one geodesic between pursuer and prey, does the optimal pursuit path switch
    to the complementary geodesic? What is the measure of such switching events?

  \item \textbf{Stealth vs.\ efficiency.} What is the Pareto frontier between
    capture time and geodesic curvature expenditure? Does the stealth constraint
    $\kappa_g \leq \kappa_{\max}$ admit an optimal control law in closed form
    for any family of metric perturbations?

  \item \textbf{Topological sufficiency of packs.} What is the minimum pack
    size $n^*$ on $\mathbb{T}^2$ that eliminates all shy coupling possibilities
    for the prey? How does $n^*$ change when the metric is dynamically
    perturbed?

  \item \textbf{Emergence and stability of neural strategies.} Does
    evolutionary training converge to a strategy that is Lyapunov-stable under
    metric perturbation, noise, and delay? Can the learned policy be certified
    by a Lyapunov function found via the perturbation theory of the
    closed-loop vector field?

  \item \textbf{Computational faithfulness.} Can the Vulkan pipeline evaluate
    the perturbed geodesic distance, the corrected SDE drift, and the curvature
    constraint projection at interactive frame rates, with a timestep $\Delta t$
    guaranteed by the Lipschitz constant of the metric field and the maximum
    curvature of active agent paths?
\end{enumerate}

% ============================================================
\section*{Dependency Graph}
% ============================================================

The following diagram encodes the strict logical dependencies of this project.
No node may be treated as understood until all nodes feeding into it are
mastered. Arrows denote ``is required by.''

\vspace{1em}

\begin{center}
\begin{tikzpicture}[
  node distance = 0.85cm and 1.4cm,
  every node/.style = {
    rectangle, rounded corners=4pt,
    minimum width=2.9cm, minimum height=0.65cm,
    font=\small\bfseries, align=center
  },
  arr/.style = {-{Stealth[length=5pt]}, thick, color=deepblue!70},
  stA/.style = {fill=stageA!18, draw=stageA, text=stageA},
  stB/.style = {fill=stageB!18, draw=stageB, text=stageB},
  stC/.style = {fill=stageC!18, draw=stageC, text=stageC},
  stD/.style = {fill=stageD!18, draw=stageD, text=stageD},
  stE/.style = {fill=purple!12, draw=purple!60!black, text=purple!70!black},
]

%% Column 1 -- Foundations
\node[stA] (ra)   {Real Analysis};
\node[stA, below=of ra]  (ms)   {Metric Spaces\\[-2pt]\&\ Topology};
\node[stA, below=of ms]  (prob) {Measure-Theoretic\\[-2pt]Probability};
\node[stA, below=of prob](na)   {Geometric Numerical\\[-2pt]Integration};

%% Column 2 -- Geometry
\node[stB, right=2.2cm of ra]  (ism)  {Smooth Manifolds\\[-2pt](Lee ISM)};
\node[stB, below=of ism] (irm)  {Riemannian\\[-2pt]Geometry (IRM)};
\node[stB, below=of irm] (pt)   {Perturbation\\[-2pt]Theory};
\node[stB, below=of pt]  (algt) {Algebraic Topology\\[-2pt]$\pi_1$, Homology};

%% Column 3 -- Dynamics
\node[stC, right=2.2cm of ism] (ode)  {ODEs on\\[-2pt]Manifolds};
\node[stC, below=of ode]  (dde)  {Delay Differential\\[-2pt]Equations};
\node[stC, below=of dde]  (ito)  {It\^{o} Calculus\\[-2pt]\&\ SDEs};
\node[stC, below=of ito]  (ocp)  {Optimal Control\\[-2pt]on Manifolds};

%% Column 4 -- Synthesis
\node[stD, right=2.2cm of ode]  (sde2) {SDEs on\\[-2pt]Manifolds};
\node[stD, below=of sde2] (sdde) {Stochastic DDEs\\[-2pt]on $\mathbb{T}^2$};
\node[stD, below=of sdde] (lyap) {Lyapunov\\[-2pt]Stability};
\node[stD, below=of lyap] (pack) {Pack Strategy\\[-2pt]\&\ Shy Couplings};

%% Column 5 -- Intelligence + Render
\node[stE, right=2.2cm of sde2] (stlth){Stealth\\[-2pt]Control};
\node[stE, below=of stlth] (nn)   {Neural Policy\\[-2pt]$\pi_\theta : \mathcal{S}\to T\mathbb{T}^2$};
\node[stE, below=of nn]    (evo)  {Evolutionary\\[-2pt]Training};
\node[stE, below=of evo]   (vk)   {Vulkan\\[-2pt]Simulation};

%% ---- Arrows: Column 1 internal ----
\draw[arr] (ra)   -- (ms);
\draw[arr] (ms)   -- (prob);
\draw[arr] (prob) -- (na);

%% ---- Arrows: Col 1 -> Col 2 ----
\draw[arr] (ra)   -- (ism);
\draw[arr] (ms)   -- (ism);

%% ---- Arrows: Col 2 internal ----
\draw[arr] (ism)  -- (irm);
\draw[arr] (irm)  -- (pt);
\draw[arr] (ms)   -- (algt);
\draw[arr] (ism)  -- (algt);

%% ---- Arrows: Col 2 -> Col 3 ----
\draw[arr] (ism)  -- (ode);
\draw[arr] (irm)  -- (ode);
\draw[arr] (pt)   -- (dde);
\draw[arr] (irm)  -- (ocp);
\draw[arr] (pt)   -- (ocp);

%% ---- Arrows: Col 1 -> Col 3 ----
\draw[arr] (prob) -- (ito);
\draw[arr] (na)   -- (ito);

%% ---- Arrows: Col 3 internal ----
\draw[arr] (ode)  -- (dde);
\draw[arr] (dde)  -- (ito);
\draw[arr] (ito)  -- (ocp);

%% ---- Arrows: Col 2+3 -> Col 4 ----
\draw[arr] (irm)  -- (sde2);
\draw[arr] (ito)  -- (sde2);
\draw[arr] (sde2) -- (sdde);
\draw[arr] (dde)  -- (sdde);
\draw[arr] (pt)   -- (sdde);
\draw[arr] (sdde) -- (lyap);
\draw[arr] (pt)   -- (lyap);
\draw[arr] (lyap) -- (pack);
\draw[arr] (algt) -- (pack);

%% ---- Arrows: Col 4 -> Col 5 ----
\draw[arr] (ocp)  -- (stlth);
\draw[arr] (lyap) -- (stlth);
\draw[arr] (sde2) -- (nn);
\draw[arr] (stlth)-- (nn);
\draw[arr] (pack) -- (nn);
\draw[arr] (nn)   -- (evo);
\draw[arr] (lyap) -- (evo);
\draw[arr] (stlth)-- (vk);
\draw[arr] (evo)  -- (vk);
\draw[arr] (pack) -- (vk);
\draw[arr] (na)   -- (vk);

\end{tikzpicture}
\end{center}

\vspace{0.8em}

\begin{center}
\small
\begin{tabular}{ll@{\hspace{1.6em}}ll@{\hspace{1.6em}}ll@{\hspace{1.6em}}ll@{\hspace{1.6em}}ll}
  \textcolor{stageA}{$\blacksquare$} & \textit{Foundations}
  & \textcolor{stageB}{$\blacksquare$} & \textit{Geometry}
  & \textcolor{stageC}{$\blacksquare$} & \textit{Dynamics}
  & \textcolor{stageD}{$\blacksquare$} & \textit{Synthesis}
  & \textcolor{purple!70!black}{$\blacksquare$} & \textit{Intelligence \& Render}
\end{tabular}
\end{center}

\vspace{1.2em}

\begin{remark}
  Perturbation theory (Stage 2) is not a terminal node. It is a \emph{vertical
  thread} that re-enters at every subsequent stage: as linearized curvature in
  Riemannian geometry, as regular and singular perturbation of ODEs and DDEs,
  as the $\varepsilon$-expansion of the stochastic heat kernel, as the
  robustness theory for Lyapunov functions, and as the metric perturbation
  budget that governs the Vulkan adaptive timestep. Likewise, algebraic
  topology enters not only in the pack strategy but in the very definition of
  the torus as a quotient---it is load-bearing from Stage 1 onward.
\end{remark}

\begin{remark}
  The C\texttt{++} engine architecture mirrors the mathematical hierarchy
  exactly. Every agent implements a drift vector field and diffusion morphism.
  Every environmental effect is a section of an appropriate bundle. The
  compositor combines them via fiberwise bundle operations. Geometric
  correctness is enforced by the interface, not by convention. The code is
  the mathematics.
\end{remark}

\vspace{0.5em}
\noindent\rule{\linewidth}{0.4pt}
\vspace{0.3em}
{\small\color{gray}\itshape
This document is the front matter of an ongoing solo research repository.
It will be revised as understanding deepens. The mathematics does not wait
for permission to be beautiful.}
